% Part: sets-functions-relations
% Chapter: functions
% Section: partial-functions

\documentclass[../../../include/open-logic-section]{subfiles}


\begin{document}

\olfileid[ar]{sfr}{fun}{par}

\olsection{الدوال الجزئية}

\begin{explain}
قد ينفع إسقاط شرط من تعريف الدالة، فلا نوجب وجود مخرج لكل مدخل
ممكن. والتعيينات التي تجوز عند هذا التخفيف تسمى \emph{دوالًا جزئية}.
\end{explain}

\begin{defn}
\emph{الدالة الجزئية} ${\OLId{latin-ordinary}{f}} \colon {\OLId{latin-looped}{A}} \pto {\OLId{latin-looped}{B}}$ تعيين يجعل لكل
!!{element} من~${\OLId{latin-looped}{A}}$ ما لا يزيد على !!{element} واحد من~${\OLId{latin-looped}{B}}$.
فإن عيّنت ${\OLId{latin-ordinary}{f}}$ عنصرًا من~${\OLId{latin-looped}{B}}$ للمدخل ${\OLId{latin-ordinary}{x}} \in {\OLId{latin-looped}{A}}$، كانت ${\OLId{latin-ordinary}{f}}({\OLId{latin-ordinary}{x}})$
\emph{معرّفة}، وإلا كانت \emph{غير معرّفة}. ونكتب في الحال الأولى،
إذا تعينت ${\OLId{latin-ordinary}{f}}({\OLId{latin-ordinary}{x}})$، الرمز ${\OLId{latin-ordinary}{f}}({\OLId{latin-ordinary}{x}}) \fdefined$، وفي الأخرى ${\OLId{latin-ordinary}{f}}({\OLId{latin-ordinary}{x}}) \fundefined$.
و\emph{مجال} الجزئية~${\OLId{latin-ordinary}{f}}$ ما عُرّفت عليه من~${\OLId{latin-looped}{A}}$، أي
$\dom{{\OLId{latin-ordinary}{f}}} = \Setabs{{\OLId{latin-ordinary}{x}} \in {\OLId{latin-looped}{A}}}{{\OLId{latin-ordinary}{f}}({\OLId{latin-ordinary}{x}}) \fdefined}$.
\end{defn}

\begin{ex}
كل دالة ${\OLId{latin-ordinary}{f}}\colon {\OLId{latin-looped}{A}} \to {\OLId{latin-looped}{B}}$ جزئية أيضًا. والجزئية التي تعرّف على~${\OLId{latin-looped}{A}}$
كلها، وهي ما سميناه إلى الآن دالة فحسب، تسمى دالة \emph{كلية}.
\end{ex}

\begin{ex}
خذ الجزئية ${\OLId{latin-ordinary}{f}} \colon \Real \pto \Real$ المحددة بـ${\OLId{latin-ordinary}{f}}({\OLId{latin-ordinary}{x}}) = {\OLMathNumeral{1}}/{\OLId{latin-ordinary}{x}}$؛
فهي لا تعرّف عند ${\OLId{latin-ordinary}{x}} = {\OLMathNumeral{0}}$، وتعرّف في كل ما سواه.
\end{ex}

\begin{prob}
خذ ${\OLId{latin-ordinary}{f}}\colon {\OLId{latin-looped}{A}} \pto {\OLId{latin-looped}{B}}$، واجعل ${\OLId{latin-ordinary}{g}}\colon {\OLId{latin-looped}{B}} \pto {\OLId{latin-looped}{A}}$ جزئية على الوجه
الآتي: لكل ${\OLId{latin-ordinary}{y}} \in {\OLId{latin-looped}{B}}$، إن وجد ${\OLId{latin-ordinary}{x}} \in {\OLId{latin-looped}{A}}$ وحيد يحقق ${\OLId{latin-ordinary}{f}}({\OLId{latin-ordinary}{x}}) = {\OLId{latin-ordinary}{y}}$
ضع ${\OLId{latin-ordinary}{g}}({\OLId{latin-ordinary}{y}}) = {\OLId{latin-ordinary}{x}}$، وإلا فـ${\OLId{latin-ordinary}{g}}({\OLId{latin-ordinary}{y}}) \fundefined$. أثبت أنه إذا كانت
${\OLId{latin-ordinary}{f}}$ متباينة، كان ${\OLId{latin-ordinary}{g}}({\OLId{latin-ordinary}{f}}({\OLId{latin-ordinary}{x}})) = {\OLId{latin-ordinary}{x}}$ لكل ${\OLId{latin-ordinary}{x}} \in \dom{{\OLId{latin-ordinary}{f}}}$،
و${\OLId{latin-ordinary}{f}}({\OLId{latin-ordinary}{g}}({\OLId{latin-ordinary}{y}})) = {\OLId{latin-ordinary}{y}}$ لكل ${\OLId{latin-ordinary}{y}} \in \ran{{\OLId{latin-ordinary}{f}}}$.
\end{prob}

\begin{defn}[الرسم البياني لدالة جزئية]
للدالة الجزئية ${\OLId{latin-ordinary}{f}}\colon {\OLId{latin-looped}{A}} \pto {\OLId{latin-looped}{B}}$ \emph{رسم بياني} للدالة~${\OLId{latin-ordinary}{f}}$، وهو العلاقة
${\OLId{latin-looped}{R}}_{\OLId{latin-ordinary}{f}} \subseteq {\OLId{latin-looped}{A}} \times {\OLId{latin-looped}{B}}$ ذات التعريف
\[
{\OLId{latin-looped}{R}}_{\OLId{latin-ordinary}{f}} = \Setabs{\tuple{{\OLId{latin-ordinary}{x}},{\OLId{latin-ordinary}{y}}}}{{\OLId{latin-ordinary}{f}}({\OLId{latin-ordinary}{x}}) = {\OLId{latin-ordinary}{y}}}.
\]
\end{defn}

\begin{prop}
لتكن ${\OLId{latin-looped}{R}} \subseteq {\OLId{latin-looped}{A}} \times {\OLId{latin-looped}{B}}$ بحيث لا يجتمع $Rxy$ و$Rxy'$ إلا مع
${\OLId{latin-ordinary}{y}} = {\OLId{latin-ordinary}{y}}'$. فتكون ${\OLId{latin-looped}{R}}$ رسم الجزئية ${\OLId{latin-ordinary}{f}}\colon {\OLId{latin-looped}{A}} \pto {\OLId{latin-looped}{B}}$ التي نعرّفها:
إن وجد ${\OLId{latin-ordinary}{y}}$ مع $Rxy$ وضعنا ${\OLId{latin-ordinary}{f}}({\OLId{latin-ordinary}{x}}) = {\OLId{latin-ordinary}{y}}$، وإلا فـ${\OLId{latin-ordinary}{f}}({\OLId{latin-ordinary}{x}}) \fundefined$.
فإن كانت ${\OLId{latin-looped}{R}}$ \emph{تسلسلية} أيضًا، أي لكل ${\OLId{latin-ordinary}{x}} \in {\OLId{latin-looped}{A}}$ يوجد
${\OLId{latin-ordinary}{y}} \in {\OLId{latin-looped}{B}}$ يحقق $Rxy$، كانت ${\OLId{latin-ordinary}{f}}$ كلية.
\end{prop}

\begin{proof}
افرض وجود ${\OLId{latin-ordinary}{y}}$ بحيث $Rxy$. فلا يجوز وجود ${\OLId{latin-ordinary}{y}}' \neq {\OLId{latin-ordinary}{y}}$ مع $Rxy'$،
وإلا خولف شرط~${\OLId{latin-looped}{R}}$. فمتى وجد ${\OLId{latin-ordinary}{y}}$ يحقق $Rxy$ كان هذا ${\OLId{latin-ordinary}{y}}$ وحيدًا،
فيصح تعريف ${\OLId{latin-ordinary}{f}}$. ومنه يظهر ${\OLId{latin-looped}{R}}_{\OLId{latin-ordinary}{f}} = {\OLId{latin-looped}{R}}$، وتكون ${\OLId{latin-ordinary}{f}}$ كلية متى
كانت~${\OLId{latin-looped}{R}}$ تسلسلية.
\end{proof}

\end{document}
